\appendix
\section{程序代码与运行说明}

完整可运行程序位于提交包 \texttt{code/solution.py}。程序依赖 Python、NumPy、pandas、Matplotlib 与 scikit-learn；解压提交包后，应在包含 \texttt{附件数据} 目录的项目根目录执行：
\begin{verbatim}
python code/solution.py
\end{verbatim}
程序从 \texttt{附件数据} 读取六个原始工作簿，在 \texttt{output/data} 写入结构化结果、约束审计、灵敏度和调度明细，在 \texttt{output/figures} 写入图表，并在项目根目录生成题目要求的结果工作簿。若输入缺失、字段非法、超出运行合同或硬约束审计失败，程序将显式退出而不是保留残缺正式结果。

\subsection{核心算法流程}
\begin{enumerate}
\item \texttt{read\_inputs} 读取附件并执行字段、范围及唯一性检查；
\item \texttt{forecast} 对区域—任务类型需求序列进行稳健回归预测；
\item \texttt{greedy\_schedule} 按实时属性、截止时间和资源需求排序，在完整合法域中构造确定性任务方案；
\item \texttt{build\_constraint\_audit} 逐任务、逐区域、逐时段复核唯一执行、时间窗、时延及瞬时 GPU、IT 与设施容量；
\item \texttt{no\_storage\_energy}、\texttt{build\_storage\_policy} 分别计算无储能与含储能能源流，并审计守恒、SOC、互斥和终端状态；
\item \texttt{run\_q4\_scenarios} 从同一联合基线独立复算各情景；
\item \texttt{main} 汇总结果并原子化写出正式表格、JSON、图表清单和结果工作簿。
\end{enumerate}

上述流程中的候选搜索深度受正文所列限制：程序未执行理论契约中的有限回溯、大规模邻域搜索或碳定向搜索，输出只证明已构造候选通过声明的确定性审计。
